\documentclass{beamer}
\usetheme{Madrid}

\title{CryptoPulse: Multimodal Sentiment \& Price Analysis}
\author{Alice Zheng, Bob Chen, Charlie Davis, Diana Lim}
\institute{SUTD}
\date{\today}

\begin{document}

\frame{\titlepage}

\begin{frame}
\frametitle{Team & Problem Statement}
\begin{itemize}
    \item \textbf{Problem:} Cryptocurrency volatility is hard to predict using price data alone. Public sentiment plays a crucial role.
    \item \textbf{Goal:} Build a multimodal model predicting Bitcoin price direction using:
    \begin{enumerate}
        \item Historical Price/Volume Data (Quantitative)
        \item News & Social Media Sentiment (Qualitative)
    \end{enumerate}
\end{itemize}
\end{frame}

\begin{frame}
\frametitle{Data Sources}
\begin{itemize}
    \item \textbf{Market Data:} Yahoo Finance API (BTC-USD, ETH-USD). Daily OHLCV from 2020-2024.
    \item \textbf{Sentiment Data:}
    \begin{itemize}
        \item NewsAPI / CryptoPanic (Headlines)
        \item Reddit (r/Bitcoin, r/CryptoCurrency)
    \end{itemize}
\end{itemize}
\end{frame}

\begin{frame}
\frametitle{Preliminary Analysis (Week 8)}
\begin{itemize}
    \item **Naive Model Implementation:** We implemented a persistence model ($P_{t+1} = P_t$) as a baseline.
    \item **Results:**
    \begin{itemize}
        \item MSE: 0.05
        \item The naive model lags behind trends, highlighting the need for predictive complexity.
    \end{itemize}
    \item **Visualization:** Plots showing Price History and Moving Averages have been generated (see Notebook).
\end{itemize}
\end{frame}

\begin{frame}
\frametitle{Proposed Methodology}
\begin{itemize}
    \item \textbf{Preprocessing:} VADER/TextBlob for sentiment scores. MinMax scaling for prices.
    \item \textbf{Model:} LSTM (Long Short-Term Memory) network taking a merged vector of [Price, Volume, Sentiment Score] at time $t$.
    \item \textbf{Output:} Predict Closing Price at $t+1$.
\end{itemize}
\end{frame}

\end{document}
